% diagrams/tp3/training_loop.tex : une itération de la boucle d'entraînement PyTorch
\begin{tikzpicture}[
    font=\sffamily\small, >=Stealth,
    blk/.style={draw=ECE, thick, fill=ECE!10, rounded corners=2pt, minimum height=1cm, align=center, inner sep=4pt},
    code/.style={font=\ttfamily\scriptsize, text=black!75}
]
\node[blk] (batch) {batch\\$(x, y)$};
\node[blk, right=0.9cm of batch] (model) {modèle\\\texttt{model(x)}};
\node[blk, right=0.9cm of model] (loss) {loss\\\texttt{loss\_fn(logits, y)}};
\node[blk, right=0.9cm of loss] (back) {gradients\\\texttt{loss.backward()}};
\node[blk, fill=gamboge!20, draw=gamboge, right=0.9cm of back] (step) {mise à jour\\\texttt{optimizer.step()}};
\draw[->, thick] (batch) -- (model);
\draw[->, thick] (model) -- node[above, font=\sffamily\scriptsize] {logits} (loss);
\draw[->, thick] (loss) -- (back);
\draw[->, thick] (back) -- (step);
\draw[->, thick, dashed] (step.south) -- ++(0,-0.6) -| node[pos=0.25, below, font=\sffamily\scriptsize]
    {poids modifiés ; \texttt{optimizer.zero\_grad()} avant le batch suivant} (model.south);
\end{tikzpicture}
